\section*{Problemas}

\subsection*{Enunciados de los problemas}

% Recordar
\begin{problema}
	\label{prob:9556f7f}
	Un modelo de regresión con una variable muda \texttt{Sex\_Male} (1 si el cliente es hombre, 0 si es mujer) ajustado sobre el Ingreso Mensual reporta:
	\begin{align}
		\hat{\texttt{Total\_Spend}} = 3570.91 + 0.146\,\texttt{Monthly\_Income} + 231.41\,\texttt{Sex\_Male}.
	\end{align}
	Interprete el coeficiente de \texttt{Sex\_Male} y prediga el gasto total para un cliente hombre y una clienta mujer, ambos con un ingreso mensual de $\$10{,}000$.
\end{problema}

% Comprender
\begin{problema}
	\label{prob:f520b9b}
	Un modelo incluye una interacción entre una variable muda $D$ y un predictor continuo $X$: $\hat Y=\hat\beta_0+\hat\beta_1X+\hat\beta_2D+\hat\beta_3(X\times D)$. Derive las dos ecuaciones reducidas del modelo para $D=0$ y $D=1$ por separado, e interprete qué representa específicamente el coeficiente $\hat\beta_3$ en términos de la diferencia de pendientes entre los dos grupos.
\end{problema}

% Aplicar
\begin{problema}
	\label{prob:5652f5a}
	Un analista compara dos modelos anidados sobre $n=150$ observaciones: el Modelo 1 (restringido, sin variables muda de ciudad) tiene $R_1^2=0.68$ con $p_1=1$ predictor; el Modelo 2 (completo, con las 2 variables muda de \texttt{City Tier}) tiene $R_2^2=0.74$ con $p_2=3$ predictores. Realice la Prueba $F$ Parcial ($H_0:$ ambos coeficientes de las variables muda son cero) al $\alpha=0.05$ ($F_{0.05,2,146}\approx3.06$), usando $F_{\text{parcial}}=\dfrac{(R_2^2-R_1^2)/(p_2-p_1)}{(1-R_2^2)/(n-p_2-1)}$.
\end{problema}

% Analizar
\begin{problema}
	\label{prob:43c3e82}
	\textbf{Demostración formal de la Trampa de la Variable Muda:} sea $\mathbf{X}$ la matriz de diseño de un modelo con intercepto y las $k$ variables muda completas ($D_1,\ldots,D_k$) para una variable categórica de $k$ niveles, sin omitir ningún nivel de referencia. Demuestre algebraicamente que las columnas de $\mathbf{X}$ satisfacen una dependencia lineal exacta ($\sum_{j=1}^k D_j = \mathbf{1}$, el vector de unos, para toda observación), y concluya formalmente por qué esto implica que $\mathbf{X}^T\mathbf{X}$ es exactamente singular y la Ecuación Normal no tiene solución única.
\end{problema}

% Evaluar
\begin{problema}
	\label{prob:a61b69d}
	\textbf{Equivalencia formal entre ANOVA de un factor y regresión con variables muda:} considere el modelo de ANOVA de un factor con $k$ tratamientos, $Y_{ij}=\mu+\tau_i+\epsilon_{ij}$, y su reformulación como regresión múltiple con $(k-1)$ variables muda: $Y=\beta_0+\beta_1D_1+\cdots+\beta_{k-1}D_{k-1}+\epsilon$, donde el tratamiento $k$ es el nivel de referencia. Evalúe y demuestre que: (a) $\beta_0=\mu_k$, (b) $\beta_j=\mu_j-\mu_k$ para $j=1,\ldots,k-1$, y (c) la Prueba $F$ Parcial de significancia conjunta de los $(k-1)$ coeficientes de las variables muda es algebraicamente idéntica a la Prueba $F$ global del ANOVA de un factor.
\end{problema}

% Crear
\begin{problema}
	\label{prob:a9bca82}
	Diseñe un problema original con una variable muda de dos niveles (distinto de sexo/región/ciudad ya vistos): un modelo predice el bono anual de desempeño ($Y$, en miles de dólares) en función de los años de experiencia ($X$) y el departamento del empleado, con Marketing como nivel de referencia y $D_{\text{Ing}}=1$ si el empleado pertenece a Ingeniería: $\hat Y = 5.0+0.8X+2.5D_{\text{Ing}}$. Interprete el coeficiente de $D_{\text{Ing}}$ y prediga el bono para un empleado de Ingeniería con $6$ años de experiencia y para uno de Marketing con la misma experiencia.
\end{problema}

\subsection*{Soluciones detalladas}

% Recordar
\begin{solproblema}[prob:9556f7f]
	El coeficiente $231.41$ indica que, manteniendo el ingreso mensual constante, un cliente hombre gasta en promedio $\$231.41$ más que una clienta mujer (el nivel de referencia). Para ingreso $=\$10{,}000$:
	\begin{align}
		\text{Hombre:}\ \hat Y = 3570.91+0.146(10000)+231.41(1) = 5262.32, \qquad \text{Mujer:}\ \hat Y = 3570.91+0.146(10000) = 5030.91.
	\end{align}
	La diferencia es exactamente $\$231.41$, el coeficiente de la variable muda.
\end{solproblema}

% Comprender
\begin{solproblema}[prob:f520b9b]
	Para $D=0$: $\hat Y=\hat\beta_0+\hat\beta_1X$ (la ecuación se reduce eliminando los términos con $D$). Para $D=1$: $\hat Y=(\hat\beta_0+\hat\beta_2)+(\hat\beta_1+\hat\beta_3)X$. La pendiente del grupo $D=0$ es $\hat\beta_1$; la pendiente del grupo $D=1$ es $\hat\beta_1+\hat\beta_3$. Por tanto, $\hat\beta_3$ representa exactamente la \textbf{diferencia de pendientes} entre los dos grupos: si $\hat\beta_3=0$, ambos grupos comparten la misma pendiente (rectas paralelas); si $\hat\beta_3\neq0$, el efecto de $X$ sobre $Y$ es distinto según el grupo (interacción genuina).
\end{solproblema}

% Aplicar
\begin{solproblema}[prob:5652f5a]
	\begin{align}
		F_{\text{parcial}} = \frac{(0.74-0.68)/(3-1)}{(1-0.74)/(150-3-1)} = \frac{0.06/2}{0.26/146} = \frac{0.03}{0.001781} \approx 16.85.
	\end{align}
	Como $16.85>3.06$, \textbf{se rechaza $H_0$}: las variables muda de \texttt{City Tier} son conjuntamente significativas; incluir la categoría de ciudad mejora significativamente el modelo.
\end{solproblema}

% Analizar
\begin{solproblema}[prob:43c3e82]
	Por construcción, cada variable muda $D_j$ satisface $D_j=1$ si y solo si la observación pertenece al nivel $j$, y $D_j=0$ en caso contrario. Como cada observación pertenece a \textbf{exactamente un} nivel (mutuamente excluyentes y colectivamente exhaustivos), para cada fila $i$ exactamente una $D_{ij}$ vale $1$:
	\begin{align}
		\sum_{j=1}^k D_{ij} = 1 \quad \text{para toda observación } i \implies \sum_{j=1}^k \mathbf{D}_j = \mathbf{1}_n,
	\end{align}
	donde $\mathbf{1}_n$ es la columna del intercepto. Esta es una dependencia lineal exacta entre las $(k+1)$ columnas de $\mathbf{X}=[\mathbf{1}_n\ \mathbf{D}_1\ \cdots\ \mathbf{D}_k]$: la columna del intercepto es idéntica a la suma de las $k$ columnas muda.

	Como existe esta dependencia, $\text{rango}(\mathbf{X})\leq k < k+1$. En consecuencia, $\mathbf{X}^T\mathbf{X}$ (de dimensión $(k+1)\times(k+1)$) también tiene rango $\leq k$, por lo que es \textbf{singular} (determinante cero, no invertible). La Ecuación Normal $\mathbf{X}^T\mathbf{X}\hat{\boldsymbol\beta}=\mathbf{X}^T\mathbf{Y}$ no tiene solución única: existen infinitas soluciones $\hat{\boldsymbol\beta}$ que producen el mismo ajuste, el mismo fenómeno que la colinealidad perfecta.
\end{solproblema}

% Evaluar
\begin{solproblema}[prob:a61b69d]
	\textbf{(a) y (b):} para una observación del tratamiento de referencia $k$, todas las variables muda son cero, así que el modelo se reduce a $E[Y]=\beta_0$. Bajo ANOVA, $E[Y_{kj}]=\mu+\tau_k=\mu_k$. Igualando: $\beta_0=\mu_k$. Para el tratamiento $j<k$: $D_j=1$ y las demás son cero, así que $E[Y]=\beta_0+\beta_j$. Bajo ANOVA, $E[Y_{ij}]=\mu_j$. Igualando y usando $\beta_0=\mu_k$: $\beta_j=\mu_j-\mu_k$.

	\textbf{(c):} la hipótesis $H_0:\beta_1=\cdots=\beta_{k-1}=0$ equivale, por las identidades anteriores, exactamente a $H_0:\mu_1=\cdots=\mu_{k-1}=\mu_k$ (todas las medias iguales) —precisamente la hipótesis nula del ANOVA, $H_0:\tau_1=\cdots=\tau_k=0$. Además, la SCR atribuible a las $(k-1)$ variables muda es algebraicamente idéntica a la SCTR del ANOVA (ambas miden la dispersión de las $k$ medias de grupo respecto a la media global, ponderada por tamaño de grupo), y los grados de libertad coinciden ($k-1$ y $N-k$ en ambos casos). Por tanto, $F_{\text{parcial}}=\text{SCR}/(k-1)$ sobre $\text{SCE}/(N-k)$ de la regresión es el \textbf{mismo número exacto} que $F=\text{CMTR}/\text{CME}$ del ANOVA: dos formulaciones matemáticas distintas del mismo procedimiento inferencial.
\end{solproblema}

% Crear
\begin{solproblema}[prob:a9bca82]
	El coeficiente $\hat\beta_2=2.5$ indica que, manteniendo constantes los años de experiencia, un empleado de Ingeniería recibe en promedio $\$2{,}500$ más de bono anual que uno de Marketing (el nivel de referencia). Para $X=6$ años:
	\begin{align}
		\text{Ingeniería:}\ \hat Y = 5.0+0.8(6)+2.5(1) = 5.0+4.8+2.5 = 12.3 \text{ miles.}
	\end{align}
	\begin{align}
		\text{Marketing:}\ \hat Y = 5.0+0.8(6)+2.5(0) = 5.0+4.8 = 9.8 \text{ miles.}
	\end{align}
	La diferencia es $12.3-9.8=2.5$, exactamente el coeficiente de $D_{\text{Ing}}$, como era de esperarse.
\end{solproblema}
